\documentclass{beamer}
\beamertemplatenavigationsymbolsempty
\usecolortheme{beaver}
\setbeamertemplate{blocks}[rounded=true, shadow=true]
\setbeamertemplate{footline}[page number]

\usepackage[T2A]{fontenc}
\usepackage[utf8]{inputenc}
\usepackage[english,russian]{babel}
\usepackage{amssymb,amsfonts,amsmath,mathtools,mathtext}

\newcommand{\cO}{\mathcal{O}}

\title[\hbox to 56mm{Алгоритм Frank--Wolfe в машинном обучении}]{Алгоритм Frank--Wolfe\\ в задачах машинного обучения}
\author[И.\,Н.~Игнашин]{Игорь Николаевич Игнашин\\
\small Научный руководитель: к.\,ф.-м.\,н. А.\,О.~Грабовой}
\institute{Кафедра интеллектуальных систем ФПМИ МФТИ\\
Специализация: Интеллектуальный анализ данных}
\date{Москва, 2026}

\begin{document}

% =====================================================================
% 1. Титул
% =====================================================================
\begin{frame}
\thispagestyle{empty}
\maketitle
\end{frame}

% =====================================================================
% 2. Положения, выносимые на защиту
% =====================================================================
\begin{frame}{Положения, выносимые на защиту}
    \footnotesize
    \begin{enumerate}
        \item Систематизированы модификации Frank--Wolfe по трем направлениям: использование предыдущих направлений, стохастические LMO на примере SOFW и расширение на децентрализованную постановку.
        \item Для модели Бекмана предложены NFW как обобщение CFW/BFW на $N$ прошлых направлений и WFFW как модификация FFW с экспоненциальным сглаживанием.
        \item CFW и NFW перенесены на ML-задачи без явного гессиана: кривизна аппроксимируется разностями градиентов.
        \item Для CFW доказана оценка по минимальному зазору Frank--Wolfe того же порядка, что у базового FW.
        \item Предложен SOFW как частный случай BCFW для транспортной задачи с блоками по OD-парам; доказана оценка ожидаемой сходимости для одиночного и батчевого выбора блоков.
        \item Предложен DBCFW: децентрализованный блочно-координатный FW с консенсусом и отслеживанием градиента; получена оценка сходимости порядка $O(1/t)$.
    \end{enumerate}
\end{frame}

% =====================================================================
% 3. Что сделано
% =====================================================================
\begin{frame}{Что сделано}
    \small
    \begin{itemize}
        \item \textbf{Методы.} Предложены NFW и WFFW для транспортной модели Бекмана, SOFW для OD-блоков транспортной задачи и DBCFW для распределенной блочной оптимизации.
        \item \textbf{Теория.} Получены оценки сходимости для CFW, SOFW и DBCFW; порядок базовых гарантий Frank--Wolfe сохраняется.
        \item \textbf{Эксперименты.} Проверка проведена на транспортных сетях, Mushrooms и MNIST; показан выигрыш памяти о направлениях и стохастического блочного оракула.
    \end{itemize}

    \vspace{5mm}
    \begin{center}
        \large \textbf{Спасибо за внимание!}

        \vspace{3mm}
        \normalsize Готов ответить на вопросы.
    \end{center}
\end{frame}

\end{document}
